\documentclass[10pt,a4paper]{amsart}
\usepackage{amsmath,amssymb,amsthm,graphicx}
% Retained minimal preamble for the Stacks Project chapter more-algebra.tex (More on Algebra).
% The chapter file pinned-original-source.tex carries only \input{preamble} on its first line, so no
% preamble text lives inside the chapter itself.  The Stacks Project uses one single shared preamble for
% every chapter; the snapshot stages it identically as sources/stacks-chow/preamble.tex and
% sources/stacks-etale-cohomology/preamble.tex.  The block below is copied mechanically, line for line,
% from sources/stacks-chow/preamble.tex lines 14-18 (the Xy-pic setup) and 154-217 (the native theorem
% declarations together with the mathematical macros).  Nothing was added, removed or redefined; the
% wrapper class and the amsmath/amssymb/amsthm/graphicx base packages are supplied by the assembler.
% Chapter: more-algebra.tex -- More on Algebra.
% For commutative diagrams we use Xy-pic
\usepackage[all]{xy}

% We use 2cell for 2-commutative diagrams.
\xyoption{2cell}

% Theorem environments.
%
\theoremstyle{plain}
\newtheorem{theorem}[subsection]{Theorem}
\newtheorem{proposition}[subsection]{Proposition}
\newtheorem{lemma}[subsection]{Lemma}

\theoremstyle{definition}
\newtheorem{definition}[subsection]{Definition}
\newtheorem{example}[subsection]{Example}
\newtheorem{exercise}[subsection]{Exercise}
\newtheorem{situation}[subsection]{Situation}

\theoremstyle{remark}
\newtheorem{remark}[subsection]{Remark}
\newtheorem{remarks}[subsection]{Remarks}

\numberwithin{equation}{subsection}

% Macros
%
\def\lim{\mathop{\mathrm{lim}}\nolimits}
\def\colim{\mathop{\mathrm{colim}}\nolimits}
\def\Spec{\mathop{\mathrm{Spec}}}
\def\Hom{\mathop{\mathrm{Hom}}\nolimits}
\def\Ext{\mathop{\mathrm{Ext}}\nolimits}
\def\SheafHom{\mathop{\mathcal{H}\!\mathit{om}}\nolimits}
\def\SheafExt{\mathop{\mathcal{E}\!\mathit{xt}}\nolimits}
\def\Sch{\mathit{Sch}}
\def\Mor{\mathop{\mathrm{Mor}}\nolimits}
\def\Ob{\mathop{\mathrm{Ob}}\nolimits}
\def\Sh{\mathop{\mathit{Sh}}\nolimits}
\def\NL{\mathop{N\!L}\nolimits}
\def\CH{\mathop{\mathrm{CH}}\nolimits}
\def\proetale{{pro\text{-}\acute{e}tale}}
\def\etale{{\acute{e}tale}}
\def\QCoh{\mathit{QCoh}}
\def\Ker{\mathop{\mathrm{Ker}}}
\def\Im{\mathop{\mathrm{Im}}}
\def\Coker{\mathop{\mathrm{Coker}}}
\def\Coim{\mathop{\mathrm{Coim}}}

% Boxtimes
%
\DeclareMathSymbol{\boxtimes}{\mathbin}{AMSa}{"02}

%
% Macros for moduli stacks/spaces
%
\def\QCohstack{\mathcal{QC}\!\mathit{oh}}
\def\Cohstack{\mathcal{C}\!\mathit{oh}}
\def\Spacesstack{\mathcal{S}\!\mathit{paces}}
\def\Quotfunctor{\mathrm{Quot}}
\def\Hilbfunctor{\mathrm{Hilb}}
\def\Curvesstack{\mathcal{C}\!\mathit{urves}}
\def\Polarizedstack{\mathcal{P}\!\mathit{olarized}}
\def\Complexesstack{\mathcal{C}\!\mathit{omplexes}}
% \Pic is the operator that assigns to X its picard group, usage \Pic(X)
% \Picardstack_{X/B} denotes the Picard stack of X over B
% \Picardfunctor_{X/B} denotes the Picard functor of X over B
\def\Pic{\mathop{\mathrm{Pic}}\nolimits}
\def\Picardstack{\mathcal{P}\!\mathit{ic}}
\def\Picardfunctor{\mathrm{Pic}}
\def\Deformationcategory{\mathcal{D}\!\mathit{ef}}

% --- appended 2026-10-12 (Chief): the project own definitions of the reference/slogan/history
% display environments, copied mechanically from the shared Stacks preamble lines 8-12
% (\usepackage{verbatim} plus the three \newenvironment lines, which the chapter preamble keeps
% verbatim from the project preamble). Nothing added, removed or redefined: these environments
% carry no mathematical content and the project itself typesets them as comments. Without these
% three lines a counted statement whose own text uses them fails with Environment slogan undefined.
\usepackage{verbatim}
\newenvironment{reference}{\comment}{\endcomment}
\newenvironment{slogan}{\comment}{\endcomment}
\newenvironment{history}{\comment}{\endcomment}
\usepackage{cleveref}
\crefname{originallemma}{original lemma}{original lemmas}
\Crefname{originallemma}{Original lemma}{Original lemmas}
\makeatletter
\newlabel{algebra-lemma-contained-in-radical}{{0AME}{source}{}{sourcepointer-1}{}}
\newlabel{algebra-lemma-contained-in-radical@cref}{{[originallemma][0][]0AME}{[1][0][]source}{}{}{}}
\makeatother
\crefname{originallemma}{original lemma}{original lemmas}
\Crefname{originallemma}{Original lemma}{Original lemmas}
\makeatletter
\newlabel{algebra-lemma-locally-nilpotent-unit}{{0AMG}{source}{}{sourcepointer-2}{}}
\newlabel{algebra-lemma-locally-nilpotent-unit@cref}{{[originallemma][0][]0AMG}{[1][0][]source}{}{}{}}
\makeatother
\crefname{originallemma}{original lemma}{original lemmas}
\Crefname{originallemma}{Original lemma}{Original lemmas}
\makeatletter
\newlabel{lemma-Rlim-pseudo-coherent-gives-pseudo-coherent}{{28336}{source}{}{sourcepointer-3}{}}
\newlabel{lemma-Rlim-pseudo-coherent-gives-pseudo-coherent@cref}{{[originallemma][0][]28336}{[1][0][]source}{}{}{}}
\makeatother
\crefname{originallemma}{original lemma}{original lemmas}
\Crefname{originallemma}{Original lemma}{Original lemmas}
\makeatletter
\newlabel{lemma-check-perfect-pointwise}{{20032}{source}{}{sourcepointer-4}{}}
\newlabel{lemma-check-perfect-pointwise@cref}{{[originallemma][0][]20032}{[1][0][]source}{}{}{}}
\makeatother
\crefname{originallemma}{original lemma}{original lemmas}
\Crefname{originallemma}{Original lemma}{Original lemmas}
\makeatletter
\newlabel{lemma-perfect}{{18889}{source}{}{sourcepointer-5}{}}
\newlabel{lemma-perfect@cref}{{[originallemma][0][]18889}{[1][0][]source}{}{}{}}
\makeatother
\title{The derived limit of an inverse system of perfect complexes along surjections with locally nilpotent kernels is perfect}
\author{The Stacks Project Authors}
\date{stacks-project:more-algebra.tex:local-snapshot-20261012; GNU FDL 1.2 or later}
\begin{document}
\maketitle
\setcounter{subsection}{1}
\noindent\textit{Editorial note.} All delivered mathematical text is the author's own, copied from the pinned Stacks Project chapter file more-algebra.tex; no statement or proof step was rewritten, reordered or removed.\par\medskip
\noindent\textit{Editorial note.} The chapter file pinned-original-source.tex carries only \textbackslash{}input\{preamble\} on its first line, so no preamble text lives inside the chapter. The retained file sources/stacks-more-algebra/retained-preamble.tex is a mechanical, line-for-line copy of the project's single shared preamble as staged in the snapshot, together with the project's own definitions of the reference/slogan/history display environments, which the project itself typesets as comments; nothing was added, removed or redefined by this card and no package was added.\par\medskip
\noindent\textit{Editorial note.} The source URL is the source-level Stacks Project URL, as required when more than one card is produced from one source; the per-claim tag is carried by claim\_key.\par\medskip
\noindent\textit{Editorial note.} The counted statement span includes the author's own \textbackslash{}begin\{reference\} block citing the source of the argument; it is delivered verbatim and produces no unresolved citation. The five references inside the counted proof span are recorded as non-load-bearing navigation pointers.\par\medskip
\noindent\textit{Allowed standard prerequisites (editorial).}\begin{itemize}
\item The pseudo-coherent version of the statement, giving pseudo-coherence and the base change isomorphisms. invoked at proof lines 28422-28424. Provides the starting point from which perfection is upgraded.
\item Elements mapping to units are units, and the resulting containment in the Jacobson radical. used at proof lines 28427-28430. Is what lets a surjection to a residue field factor through A\_1.
\item Perfect complexes have finite tor dimension, and the pointwise criterion for perfection. invoked at proof lines 28438-28443. Convert the boundedness of the residue field fibres into perfection of K.
\end{itemize}\medskip
\ifdefined\hypertarget\hypertarget{sourcepointer-1}{}\fi
\noindent\textit{Source reference.} Original lemma 0AME: focused-seeds/AGP/sources/stacks-snapshot-20261011/tags.txt, source lines 13402--13402. Cross-chapter navigation pointer into the Stacks Project; the standard fact it supplies is applied inside the counted proof span and the argument resting on it stays there.\par\medskip
\ifdefined\hypertarget\hypertarget{sourcepointer-2}{}\fi
\noindent\textit{Source reference.} Original lemma 0AMG: focused-seeds/AGP/sources/stacks-snapshot-20261011/tags.txt, source lines 13404--13404. Cross-chapter navigation pointer into the Stacks Project; the standard fact it supplies is applied inside the counted proof span and the argument resting on it stays there.\par\medskip
\ifdefined\hypertarget\hypertarget{sourcepointer-3}{}\fi
\noindent\textit{Source reference.} Original lemma 28336: sources/stacks-more-algebra/pinned-original-source.tex, source lines 28336--28336. Same-chapter navigation pointer; the statement it names is standard for the section and the argument that consumes it stays inside the counted proof span.\par\medskip
\ifdefined\hypertarget\hypertarget{sourcepointer-4}{}\fi
\noindent\textit{Source reference.} Original lemma 20032: sources/stacks-more-algebra/pinned-original-source.tex, source lines 20032--20032. Same-chapter navigation pointer; the statement it names is standard for the section and the argument that consumes it stays inside the counted proof span.\par\medskip
\ifdefined\hypertarget\hypertarget{sourcepointer-5}{}\fi
\noindent\textit{Source reference.} Original lemma 18889: sources/stacks-more-algebra/pinned-original-source.tex, source lines 18889--18889. Same-chapter navigation pointer; the statement it names is standard for the section and the argument that consumes it stays inside the counted proof span.\par\medskip
\begin{lemma}
\label{lemma-Rlim-perfect-gives-perfect}
\begin{reference}
\cite[Lemma 4.2]{Bhatt-Algebraize}
\end{reference}
Let $A = \lim A_n$ be a limit of an inverse system $(A_n)$ of rings.
Suppose given $K_n \in D(A_n)$ and maps $K_{n + 1} \to K_n$
in $D(A_{n + 1})$. Assume
\begin{enumerate}
\item the transition maps $A_{n + 1} \to A_n$ are surjective with
locally nilpotent kernels,
\item either all $K_n$ are perfect or there exists an integer $n_0$
such that $K_{n_0}$ is perfect and the kernels $A_{n + 1} \to A_n$
are nilpotent for $n \geq n_0$, and
\item the maps induce isomorphisms
$K_{n + 1} \otimes_{A_{n + 1}}^\mathbf{L} A_n \to K_n$.
\end{enumerate}
Then $K = R\lim K_n$ is a perfect object of $D(A)$
and $K \otimes_A^\mathbf{L} A_n \to K_n$ is an isomorphism for all $n$.
\end{lemma}

\begin{proof}
We already know that $K$ is pseudo-coherent and that
$K \otimes_A^\mathbf{L} A_n \to K_n$ is an isomorphism for all $n$
by Lemma \ref{lemma-Rlim-pseudo-coherent-gives-pseudo-coherent}.
Consider a surjective map $A \to \kappa$ whose kernel is
a maximal ideal $\mathfrak m$.
Any element of $A$ which maps to a unit in $A_1$ is a unit
in $A$ by Algebra, Lemma \ref{algebra-lemma-locally-nilpotent-unit}
and hence $\Ker(A \to A_1)$ is contained in the Jacobson radical of $A$ by
Algebra, Lemma \ref{algebra-lemma-contained-in-radical}.
Hence $A \to \kappa$ factors as $A \to A_1 \to \kappa$.
Hence
$$
K \otimes_A^\mathbf{L} \kappa =
K \otimes_A^\mathbf{L} A_1 \otimes_{A_1}^\mathbf{L} \kappa =
K_1 \otimes_{A_1}^\mathbf{L} \kappa
$$
Note that $K_1$ is perfect by assumption (2). Hence
$K_1$ has finite tor dimension by Lemma \ref{lemma-perfect}.
Thus there exist $a, b \in \mathbf{Z}$ such that
$H^i(K \otimes_A^\mathbf{L} \kappa) = 0$ for all $i \not \in [a, b]$.
By Lemma \ref{lemma-check-perfect-pointwise}
we conclude that $K$ is perfect.
\end{proof}

\end{document}
